\section*{Chapter \ref{ch:iv1:logic-riddles} --- Logic Riddles}
\markboth{\protect\omTocMark{}{Answers: Logic Riddles}}{\protect\omTocMark{}{Answers: Logic Riddles}}

\begin{solution}{iq:iv1:logic-riddles:1}
Two weighings for 8 balls, and 81 balls with 4 weighings.

A balance has three outcomes\index{information counting}: left heavy, right heavy, or level. So $k$ weighings separate at most $3^k$ cases. Put 3 against 3. If one side is heavy, the ball is among those 3. If level, it is among the other 2. Then one more weighing finds it among 3, by putting 1 against 1.

\iqwrong{Halving each time, as if the balance had only two outcomes: 4 against 4, then 2 against 2, then 1 against 1. That takes 3 weighings, and it never uses the level outcome.}
\end{solution}

\begin{solution}{iq:iv1:logic-riddles:2}
Light rope A at both ends and rope B at one end. When A burns out, after 30 minutes, light the other end of B. It burns out 15 minutes later, for 45 in all.

Burning from both ends halves the time, however uneven the rope is\index{invariant reasoning}, because the two flames together consume the whole 60 minutes of burning. When A is gone, rope B has 30 minutes of burning left, and the second flame halves that.

\iqwrong{Cutting a rope in half or in a quarter. The burn rate is uneven, so length says nothing about time. Only the time totals can be used.}
\end{solution}

\begin{solution}{iq:iv1:logic-riddles:3}
Start both. At 4, turn the 4-minute timer. At 7, turn the 7-minute timer. At 8, the 4-minute timer runs out, and the 7-minute timer has run for 1 minute, so turn it over and it runs for that 1 minute back. That ends at 9.

The trick is to use the sand that has already fallen\index{timers}. A timer that has run for $t$ minutes, once turned, measures exactly $t$ minutes. So you need to find a moment where such a timer has the amount you need: here 1 minute after the 8-minute mark.

\iqwrong{Trying to combine only multiples, as with $4a + 7b$, which does not include 9. You can only get 9 by using a turn at a moment when a timer is part of the way through.}
\end{solution}

\begin{solution}{iq:iv1:logic-riddles:4}
Take 1 coin from bag 1, 2 from bag 2, and so on up to 10 from bag 10. That is 55 coins, which would weigh 550 grams if all were genuine. The excess in grams is the number of the heavy bag.

One reading of a scale can carry many values, so you make each bag contribute a different amount\index{coding by weights}. The excess is $k$ when bag $k$ is the fake one, and this is different for each of the ten cases.

\iqwrong{Weighing one bag at a time, which needs up to nine uses of the scale. A scale is a number, not a yes or a no.}
\end{solution}

\begin{solution}{iq:iv1:logic-riddles:5}
Take one more A from its bottle, so that you hold 2 A and 2 B. Cut each of the four pills in half, and put one half of each into today's pile, and the other half into tomorrow's.

You cannot tell the pills apart, so you cannot select by type, but you can make the selection the same for every pill\index{symmetrising}. Today's four halves are half of 2 A and half of 2 B, which is exactly one A and one B. Tomorrow's four halves are the same dose. No pill is wasted.

\iqwrong{Throwing the pills away and starting afresh, which fails the question, or taking any two of the three, which risks a wrong dose.}
\end{solution}

\begin{solution}{iq:iv1:logic-riddles:6}
One cut is enough, at the third link (or, by symmetry, the fifth), which leaves pieces of 1, 2 and 4 links. With 2 cuts, a chain of 23 links works.

On day $d$ the worker must hold exactly $d$ links, so the pieces must be able to make every total from 1 to 7, using each piece at most once. Pieces of 1, 2 and 4 do this, as in binary\index{binary representation}. So on day 2 you take back the 1 and hand over the 2.

With 2 cuts you get two single links, which reach 2. The next piece may then be 3, reaching 5, the next 6, reaching 11, and the last 12, reaching 23. So $1 + 1 + 3 + 6 + 12 = 23$. Trying every pair of cuts confirms that no longer chain works.

\iqwrong{Three cuts, at links 2, 4 and 6, which leave seven single links to hand over one a day. This ignores change, which is the point of the puzzle.}
\end{solution}

\enlargethispage{-3\baselineskip}%
\begin{solution}{iq:iv1:logic-riddles:7}
No for the opposite corners. Yes for any two squares of opposite colours.

Colour the board in the usual way. A domino covers one black and one white square\index{colouring argument}. Opposite corners have the same colour, so the rest has 30 of one colour and 32 of the other, and dominoes cannot cover that.

For two squares of opposite colours, there is a nice proof. Draw a closed path through all 64 squares, moving between neighbours, for example a comb shape. Removing two squares of opposite colours breaks the path into two pieces, each with an even number of squares, because colours alternate along the path. Each piece is covered by dominoes laid along it. An exhaustive search on a 4 by 4 board agrees.

\iqwrong{Trying to place the dominoes one by one and getting stuck, and saying that it seems to be impossible. That proves nothing. A colour count does it in one line.}
\end{solution}

\begin{solution}{iq:iv1:logic-riddles:8}
One always says the colour they see on the other. The other always says the opposite of what they see. Exactly one of them is right.

If the hats match, the first is right. If the hats differ, the second is right\index{hat puzzles}. These are the only two cases, so one of them always wins. Each person is right half of the time, as they must be, because they know nothing about their own hat. The plan does not change that. It arranges that the right guesses never coincide.

\iqwrong{Saying this cannot be done, since each guess is a coin flip. It is true for each guess alone. A plan can correlate the guesses.}
\end{solution}

\begin{solution}{iq:iv1:logic-riddles:9}
Turn switch 1 on for ten minutes, turn it off, then turn switch 2 on, and enter. The lit bulb is switch 2, the warm unlit bulb is switch 1, and the cold unlit bulb is switch 3.

One look at an on or off state gives only two values, which is not enough for three switches\index{information counting}. A filament bulb has a further state, its temperature. You need a third signal, and heat is one you can create in advance.

\iqwrong{Saying it cannot be done, because there is one entry and only two states. Ask what else can differ about a bulb. An interviewer may want to hear you question the model, not only the arithmetic.}
\end{solution}

\begin{solution}{iq:iv1:logic-riddles:10}
Take a fruit from the box labelled ``mixed''. It really holds only the fruit that you draw. If you draw an apple, that box is apples, the box labelled ``oranges'' cannot be oranges and it cannot be apples, so it is mixed, and the box labelled ``apples'' is oranges.

All labels are wrong, so only two arrangements are possible\index{derangements}. The box labelled ``mixed'' holds pure fruit in both of them, so one draw from it tells you which arrangement you have. A draw from the box labelled ``apples'' may show an orange in both arrangements, and in that case it tells you nothing.

\iqwrong{Choosing a box at random. A fruit from a box that may be mixed does not settle what the box contains.}
\end{solution}

\begin{solution}{iq:iv1:logic-riddles:11}
The shortest time is 17 minutes. Persons 1 and 2 cross (2 minutes), 1 returns (1), persons 5 and 10 cross (10), 2 returns (2), and 1 and 2 cross (2).

The natural plan is to use the fastest person to escort everyone, which takes $10 + 1 + 5 + 1 + 2 = 19$. The better plan sends the two slowest together\index{exchange argument}, so that the 5 costs nothing extra, at the price of needing person 2 for one return. A search over all states finds a minimum of 17.

\iqwrong{Giving 19, by assuming that the fastest person should always carry the torch. Grouping the two slowest saves more than the fast escort does.}
\end{solution}

\begin{solution}{iq:iv1:logic-riddles:12}
Fill the 5, pour into the 3, leaving 2. Empty the 3, and pour in the 2. Fill the 5, and top up the 3 with 1 litre, which leaves 4. With jugs of 6 and 10 you cannot get 3.

Any amount you can make is a combination of the two sizes, so it is a multiple of their greatest common divisor\index{B\'ezout's identity}. And you can reach every multiple of it up to the larger jug. For 3 and 5, the divisor is 1. For 6 and 10, it is 2, and so only even amounts arise. A search over all states of the jugs confirms the amounts $\{0, 2, 4, 6, 8, 10\}$.

\iqwrong{Trying moves at random until the answer appears. State the rule, which also tells you when to stop.}
\end{solution}

\begin{solution}{iq:iv1:logic-riddles:13}
No. The parity of the number of black cells never changes, and it starts at 1.

A move flips 8 cells\index{parity invariant}. If $k$ of those were black, the black count changes by $8 - 2k$, which is even. So the count stays odd, and it cannot be 0. A long random run of moves on a grid confirms the count stays odd at every step.

\iqwrong{Trying to fix the one black cell by flipping its row. That makes the row black except for one cell, and the count goes to 7: odd again.}
\iqfollow{And on a 7 by 7 grid?}{A move now flips 7 cells, so the count no longer helps. Use any 2 by 2 block instead: a move flips 0 or 2 of its cells, so its parity never changes, and a block that holds the black cell stays odd. So the answer is still no, on any size of grid.}
\end{solution}

\begin{solution}{iq:iv1:logic-riddles:14}
Ninety-nine can be sure, and the last person at the back has an even chance. The one at the back says ``black'' if the number of black hats in front is odd, and ``white'' otherwise.

The plan encodes one bit, the parity of the 99 hats\index{parity coding}. The next person sees 98 of those hats. Comparing their parity with what was announced gives their own colour for certain. Each person in turn knows the announced parity and has heard all the answers from behind, and sees the rest, so each can work out their own hat. A simulation with 300 random lines finds no error among those 99.

So, on average, $99.5$ are right. You cannot do better than that person's even chance, since nobody has seen that hat.

\iqwrong{Pairing people so that each says what the other one wears, and saving about half for sure. That wastes the fact that everybody hears everything, and one bit is enough for all.}
\iqfollow{What if there are three colours?}{Use the sum of the colour numbers modulo 3, and say it as one of three words. Still 99 are sure. In general, one announcement carries one sum modulo the number of colours.}
\end{solution}

\begin{solution}{iq:iv1:logic-riddles:15}
You need 10 tasters for one round, and 7 tasters for two rounds.

One round: write each bottle's number in binary\index{binary coding}, with 10 digits, because $2^{10} = 1024 \ge 1000$. Taster $k$ sips from every bottle whose $k$-th digit is 1. The set of those affected gives the digits of the poisoned bottle. You cannot do with 9 tasters, since 9 tasters produce only 512 distinct results.

Two rounds: now each taster has three possible results: affected in round 1, affected in round 2, or never. So write the bottles in base 3, with 7 digits, as $3^6 = 729 < 1000 \le 3^7 = 2187$. A digit of 0 means that the taster drinks in round 1, a digit of 1 means round 2, and a digit of 2 means never. A taster affected in round 1 does not need to drink in round 2, because that digit is already known.

In general, with $r$ rounds, $t$ tasters handle $(r + 1)^t$ bottles.

\iqwrong{Using 1000 tasters, or a binary search, which takes 10 days. Group testing in parallel is about designing one test whose results encode many bits at once.}
\iqwhy{This is the idea behind group testing and fast fault isolation: put many things in each test, and let each one's pattern of results be its address.}
\end{solution}

\begin{solution}{iq:iv1:logic-riddles:16}
Take any 10 coins as the first group, and flip all of them. The two groups then have the same number of heads.

Say that the 10 you chose hold $h$ heads\index{invariant reasoning}. So the other 90 hold $10 - h$. After you flip your 10, they hold $10 - h$ heads as well. It does not matter how many heads you chose, which is why you do not need to see. A simulation with random arrangements gives equality each time.

\iqwrong{Saying that you need to feel the coins, or looking for a way of dividing them evenly. The number you need is the size of the group, and it comes from the total given, not from the coins' positions.}
\end{solution}

\begin{solution}{iq:iv1:logic-riddles:17}
No. The two corners have the same colour, but a knight's path over all 64 squares has 63 moves, and it must end on the opposite colour from where it started.

Every knight's move changes the colour of the square\index{colouring argument}. After 63 moves, an odd number, you are on the opposite colour. Opposite corners of the board have the same colour, because $(0 + 0)$ and $(7 + 7)$ are both even. So it is impossible.

The argument forbids only the same colour, so a finish on the other colour is allowed by this argument. It does not by itself prove such a path exists. On small boards you can test the rule by exhaustive search: on a 4 by 4 board, no tour joins opposite corners, while on a 5 by 5 board, with 24 moves, such a tour exists.

\iqwrong{Trying to draw a tour and failing. A failure to find a tour says nothing, and on larger boards you would not finish. Parity is one line.}
\end{solution}

\begin{solution}{iq:iv1:logic-riddles:18}
Use induction on $n$. An 8 by 8 board uses 21 pieces.

For $n = 1$, a 2 by 2 board with one square removed is an L piece. Now take a $2^n$ by $2^n$ board with a square removed. Split it into four quadrants\index{proof by induction}. The missing square lies in one quadrant. Place one L piece at the centre of the board so that it takes one square from each of the other three quadrants. Now each of the four quadrants has exactly one missing square, and each is a smaller board of the same type.

The count is $(64 - 1)/3 = 21$, and $4^n - 1$ is always divisible by 3. A program carrying out this construction, for each of the 64 positions of the hole, covers everything with 21 pieces of three squares.

\iqwrong{Assuming that the count alone gives the answer. That only shows that the numbers fit, and a colouring could still forbid some positions. The induction gives you a real covering.}
\end{solution}

\begin{solution}{iq:iv1:logic-riddles:19}
Not with 13, 15 and 17. With 13, 16 and 16, yes.

When two meet, one colour count goes up by 2 and two go down by 1\index{invariant modulo 3}. So the difference between any two counts changes by 0 or by $\pm 3$, and it is unchanged modulo 3. If all were one colour, two counts would be 0, with a difference of 0. In 13, 15, 17, the residues are 1, 0, 2, all different, so no two counts can ever be equal modulo 3.

With 13, 16 and 16, two counts agree. Let them meet: each meeting of green and blue reduces both, and 16 meetings turn them all red. A search over all reachable states confirms the two answers.

\iqwrong{Thinking of parity, as with so many invariants. A meeting changes the counts by 1 and 2, so parity is not preserved, and the right modulus is 3.}
\end{solution}

\begin{solution}{iq:iv1:logic-riddles:20}
It cannot be 0, since the parity of the sum never changes and it starts odd. The second game always ends in $11! - 1 = 39\,916\,799$, whatever the order.

First, $|a - b|$ has the same parity as $a + b$\index{parity invariant}, so each move keeps the parity of the total. The total starts at 55, which is odd, so the last number is odd.

Second, note that $a + b + ab = (a + 1)(b + 1) - 1$\index{invariant product}. So the product of $x + 1$ over all the numbers on the board does not change. It starts as $2 \cdot 3 \cdots 11 = 11!$, and at the end it is the last number plus one. Random orders in a simulation all confirm these two results.

\iqwrong{Trying a few orders to look for a pattern. Any invariant will be easier to find by asking what a move preserves.}
\end{solution}

\begin{solution}{iq:iv1:logic-riddles:21}
Use 148 comparisons for 100 numbers. In general it is $\lceil 3n/2 \rceil - 2$.

Compare the numbers in pairs\index{comparison counting}. Each pair uses one comparison, and then the larger of each pair competes for the maximum, and the smaller of each pair for the minimum. For $n = 100$: 50 comparisons in pairs, then 49 for the maximum, and 49 for the minimum, giving 148. A count over real arrays of various sizes confirms the formula.

Why it is hard to beat: to finish, $n - 1$ numbers must have lost once, and $n - 1$ must have won once, so $2n - 2$ units of information are needed. A comparison between two untouched numbers gives 2 units, and any later one gives at most 1 against an adversary. There are at most $n/2$ of the first type.

\iqwrong{Using $198$ comparisons, or claiming that sorting first is good. Sorting costs about $n\log_2 n$. Using a single scan for both ignores what each comparison can tell you.}
\end{solution}

\enlargethispage{-3\baselineskip}%
\begin{solution}{iq:iv1:logic-riddles:22}
The ages are 2, 2 and 9, and the house number is 13.

List the triples with product 36. There are eight: $(1, 1, 36)$, $(1, 2, 18)$, $(1, 3, 12)$, $(1, 4, 9)$, $(1, 6, 6)$, $(2, 2, 9)$, $(2, 3, 6)$, $(3, 3, 4)$. Their sums are 38, 21, 16, 14, 13, 13, 11 and 10. The census taker sees the number, and still cannot tell, so it must be a sum that appears twice\index{knowledge as information}: 13. Then ``the eldest'' rules out $(1, 6, 6)$, where the two oldest are the same age.

\iqwrong{Stopping when the sum looks ambiguous to you. The information is in the fact that the census taker could not decide, which is a statement about the taker's view, not yours.}
\end{solution}

\begin{solution}{iq:iv1:logic-riddles:23}
The hat at the front is red.

There are only 2 white hats, so if the back person saw two white hats, they would know that their own is red\index{knowledge as information}. They say they do not know, so the middle and the front are not both white. The middle person knows this. If the front hat were white, the middle person would see it and know that their own must be red. But the middle person does not know, so the front hat is not white.

The front person sees nobody, and works from what the other two said. Each ``I do not know'' is a piece of information about what they saw.

\iqwrong{Saying that the front person has no information because they see nothing. Each statement narrows the possible worlds, from seven down to those with a red hat at the front.}
\end{solution}

\begin{solution}{iq:iv1:logic-riddles:24}
Take $2^{k - 1}$ coins from stack $k$: 1, 2, 4, up to 512, for 1023 coins. The excess in grams, written in binary, lists the fake stacks.

Each subset of stacks must give a different excess\index{binary coding}. Powers of 2 do this, because every number has a unique binary expansion. The 1024 possible subsets need 1024 distinct excesses, starting from 0, so a total of at least 1023 extra grams must be possible, which means at least 1023 coins, and the powers of 2 reach this bound. An exhaustive check of all 1024 subsets agrees.

\iqwrong{Using 1, 2, 3, up to 10, as you would for a single fake stack. Then 3 might mean stack 3, or stacks 1 and 2 together.}
\end{solution}

\begin{solution}{iq:iv1:logic-riddles:25}
Ask one guard, ``Would you say yes if I asked you whether the left road is safe?'' A yes means that left is safe, in all cases.

Check it. If left is safe and the guard is truthful, the direct answer would be yes, and they say yes. If left is safe and the guard lies, their direct answer would be no, and they lie about that, so they say yes. If left is not safe, the truthful guard says no, and the liar's direct answer would be yes, which they lie about, so they say no\index{double negation}.

Another form uses the other guard: ``Would the other guard say that left is safe?'' Here the answer is always false, so a no means that left is safe. An exhaustive check of the four cases confirms both.

\iqwrong{Asking a question that identifies the liar, such as ``Is two plus two four?'' This uses your only question on something other than the road.}
\end{solution}

\begin{solution}{iq:iv1:logic-riddles:26}
Statement 99 is the only true one. With ``at least $k$'', statements 1 to 50 are true, and the rest are false.

At most one ``exactly'' statement can be true, because they contradict one another. If none were true, all 100 would be false, which would make statement 100 true: a contradiction. So exactly one is true, and 99 are false, which makes statement 99 the true one\index{self-reference}.

For ``at least $k$'', suppose that $f$ statements are false. The true ones are exactly those with $k \le f$, so there are $f$ true ones. Hence $100 - f = f$, and $f = 50$. A search over all possible numbers of false statements finds that each version has one solution.

\iqwrong{Choosing statement 100, because it is the last one. If it were true, all 100 would be false, including itself.}
\end{solution}

\begin{solution}{iq:iv1:logic-riddles:27}
Use a fixed scheme: weigh $\{1, 2, 3, 10\}$ against $\{4, 5, 6, 11\}$, then $\{1, 2, 3, 11\}$ against $\{7, 8, 9, 10\}$, and then $\{1, 4, 7, 10\}$ against $\{2, 5, 8, 12\}$. Thirteen balls cannot be done.

Each ball has a pattern over the three weighings, $+$, $-$ or $0$ for left, right or out. If the ball is heavy, that pattern is the outcome. If it is light, it is the negative of that pattern\index{ternary coding}. So all 24 cases must give different outcomes, and the pattern of each ball must not be the negative of another's. In this scheme, ball 1 is $(+, +, +)$, ball 2 is $(+, +, -)$, ball 10 is $(+, -, +)$ and ball 11 is $(-, +, 0)$. An exhaustive check confirms that the 24 outcomes are all different and none is level in all three.

Now for why 13 fails: $27 - 1 = 26$ outcomes and $13 \cdot 2 = 26$ cases seem to fit exactly. But look at the first weighing. Put $a$ balls on each pan. If it tips, there are $2a$ cases left, and two weighings have 9 outcomes, so $a \le 4$. If level, there are $13 - 2a$ balls left, hence $2(13 - 2a)$ cases, so that must be at most 9, which gives $a \ge 5$. These contradict each other.

\iqwrong{Thinking that $3^3 = 27$ settles the count. That count is a bound. You also need a first step that splits the cases evenly, and the need for equal numbers on both pans rules this out for 13.}
\end{solution}

\begin{solution}{iq:iv1:logic-riddles:28}
No for 10 by 10, and yes for 4 by 10. An $m$ by $n$ board can be covered exactly when 4 divides $m$ or $n$.

Colour the cell at $(x, y)$ with $(x + y) \bmod 4$\index{colouring argument}. Each straight brick covers one cell of each of the four colours, whichever way it lies. On a 10 by 10 board, the colour counts are 25, 26, 25 and 24, which are not equal. So no covering exists.

For the general rule, use a stronger colouring. Put the number $i^{x + y}$ on each cell, with $i^2 = -1$. A brick has four consecutive powers of $i$, and they add up to 0\index{roots of unity}. So the sum over a covered board is 0. But the sum over the board is
\[
\Bigl(\sum_{x = 0}^{m - 1} i^x\Bigr)\Bigl(\sum_{y = 0}^{n - 1} i^y\Bigr),
\]
and a sum of consecutive powers of $i$ is 0 only when its length is a multiple of 4. So 4 must divide $m$ or $n$. And if it does, laying the bricks in lines is enough. An exhaustive search for all sizes up to 7 by 7 agrees with this rule.

\iqwrong{Using the total count of cells. A 10 by 10 board has 100 cells, which is divisible by 4, and it still cannot be covered. A count is only a first check.}
\end{solution}

\begin{solution}{iq:iv1:logic-riddles:29}
For 64 numbers, 68 comparisons, against $63 + 62 = 125$ for two scans. In general it is $n + \lceil \log_2 n \rceil - 2$.

Run a knockout tournament\index{tournament method}: $n - 1$ comparisons, which find the largest. The second-largest must have lost directly to the largest, because it could not lose to anything else. So it is among the numbers that the winner beat, of which there are $\lceil \log_2 n \rceil$. Scan those for the largest with $\lceil \log_2 n \rceil - 1$ more comparisons. For $n = 64$, that is $63 + 5$. A run on random arrays gives the second-largest correctly, with a worst case of exactly 68.

Why it must be so: you must know that every number except the top two is not second. The only proof that a number is not second is that it lost to something other than the largest. Otherwise, it could yet be second. So $n - 2$ numbers need a loss to a non-largest. On top of that, the largest has to win its own comparisons, and an adversary can make it win at least $\lceil \log_2 n \rceil$ of them, none of which count toward the first group\index{adversary argument}.

\iqwrong{Sorting, or removing the largest and rescanning, for 125 comparisons. That second scan repeats comparisons whose results you already know.}
\end{solution}

\begin{solution}{iq:iv1:logic-riddles:30}
Number the people from 0 to $n - 1$, and the colours from 0 to $n - 1$. Person $k$ guesses the colour that would make the sum of all colours equal $k$ modulo $n$. Exactly one is right.

Let $S$ be the true sum, modulo $n$\index{hat puzzles}\index{modular coding}. Person $k$ sees all hats but their own, and assumes that the total is $k$. So person $k$ is right exactly when $S = k$. As $S$ takes one value from $0$ to $n - 1$, exactly one person is right.

Every person is right with chance $1/n$, as they must be, because each hat is independent of what they see. The total number of correct guesses across the group always has a mean of 1. Without a plan, this number varies, and it may be 0. The plan removes all of the variance. An exhaustive check of all arrangements for $n$ up to 5 finds exactly one right guess every time.

\iqwrong{Thinking that a plan must raise each person's chance. It cannot, so it must change only how right guesses overlap.}
\iqwhy{Diversification works the same way. You cannot change each bet's expected payoff, but you can arrange the correlation so that exactly one pays in each state.}
\end{solution}

\begin{solution}{iq:iv1:logic-riddles:31}
They can win with probability $3/4$. The plan: if you see two hats of the same colour, guess the opposite colour. If you see two different hats, pass.

There are eight outcomes. In the 6 with two of one colour and one of the other, the person with the odd hat sees two equal hats and guesses correctly, and the other two see different hats and pass. That is a win\index{Hamming code}. In the two where all three match, all three guess wrongly, which is a loss.

Each guess is still right half of the time. Across the eight outcomes, the plan makes 12 guesses: 6 are right, one in each of the winning outcomes, and 6 are wrong, all in the two losing outcomes. The plan piles the wrong guesses up into a few outcomes.

No plan beats this. Right and wrong guesses are equal in number for any plan. Each winning outcome needs at least one right guess, and each losing outcome has at most 3 wrong ones. So wins $\le 3 \times$ losses, which gives at most $6$ of 8. A search over all $81^3$ plans confirms a maximum of 6.

\iqwrong{Saying that the chance is $1/2$, because any guess has chance $1/2$. Only the group result counts, and a plan changes how right and wrong guesses overlap, not their total numbers.}
\iqfollow{And with seven people?}{A win chance of $7/8$, using the Hamming code of length 7: guess only if the hats you see, with some colour for your hat, would be a codeword, and choose the other colour.}
\end{solution}

\begin{solution}{iq:iv1:logic-riddles:32}
All 100 blue-eyed people leave together on the 100th midnight, and nobody else leaves. The visitor made a fact common knowledge, which it was not before.

Use induction on the number $k$ of blue-eyed people\index{common knowledge}. If $k = 1$, that person sees no blue eyes, and leaves on night 1. If $k = 2$, each sees one blue-eyed person. Each thinks, ``If I am not blue, the other will leave tonight.'' The other does not, so both leave on night 2. In general, each of $k$ blue-eyed people sees $k - 1$, and waits to see whether those leave on night $k - 1$. When they do not, they all leave on night $k$. A brown-eyed person sees 100 blue, and when those leave on night 100, learns only that their own eyes are not blue, which does not give their colour.

What was new? Everybody knew that someone had blue eyes. But with two blue-eyed, A does not know that B knows. With three, A does not know that B knows that C knows. With 100, the chain of ``knows that'' has to be 100 levels long, and it stops short without a public statement. The announcement makes every level true at once. A possible-worlds simulation with up to 5 blue-eyed people among 6 reproduces the days.

\iqwrong{Saying that nothing happens, since the visitor told them nothing they did not know. The first level of knowledge did not change. The higher levels, what each knows about what others know, did change.}
\iqwhy{Public news can move a market even when everybody knew it privately. If you do not know that others know, you act differently from when you do.}
\end{solution}

\begin{solution}{iq:iv1:logic-riddles:33}
One prisoner is the counter. Each of the others turns the light on the first time that they find it off, and never again. Only the counter turns it off, and adds 1 each time. When the counter reaches 99, they declare. This takes about $10\,418$ days on average, or about $28.5$ years.

It cannot fail, because the count only goes up when someone new has signalled\index{counter protocol}. The switch starts off, so there is no false count from the start. Each of the 99 non-counters turns the light on at most once.

The time is a sum of two waits for each of the 99 signals. If $100 - k$ non-counters have not yet signalled, one of them comes along after $\frac{100}{100 - k}$ days on average. Then the counter must come, which takes 100 days on average. So
\[
\sum_{k = 1}^{99} \Bigl(\frac{100}{100 - k} + 100\Bigr) = 100 H_{99} + 9900 \approx 10\,418.
\]
Nearly all of this is waiting for the counter. A simulation of the plan gives a mean within $1.5\%$ of this.

\iqwrong{Declaring after some large number of days, on the grounds that everyone must have visited by then. That can be made very likely, but the question requires certainty. A count is sure, and any time limit is not.}
\iqfollow{What if the switch's start is unknown?}{Then count to twice as much: each non-counter turns it on twice, and the counter stops at $2 \cdot 99 = 198$, which takes about twice as long. The one initial false count cannot reach the total by itself.}
\end{solution}

\begin{solution}{iq:iv1:logic-riddles:34}
Number the squares from 0 to 63. Let $P$ be the exclusive-or of the numbers of all squares that show heads. Prisoner A flips the coin at square $P \oplus t$, where $t$ is the target. Prisoner B computes the new $P$, which equals $t$. The method works exactly when the number of squares is a power of 2.

Flipping the coin on square $s$ changes $P$ to $P \oplus s$\index{parity coding}\index{exclusive-or}, whether that coin was heads or tails, because a number's own exclusive-or with itself is 0. So A wants $P \oplus s = t$, which is $s = P \oplus t$. This is always a valid square, as $P$ and $t$ both have 6 bits. So every target is available to A, in exactly one flip. A test with random boards and targets finds no failure.

What is special about 64: the plan assigns a square to each arrangement, and from any arrangement, one flip must reach each target. On the cube of arrangements with $n$ coordinates, each corner has $n$ neighbours, and all $n$ colours have to appear among them. So each colour is on exactly $2^n/n$ corners, and $n$ must divide $2^n$. So $n$ must be a power of 2. An exhaustive search for three squares finds no colouring, as this argument predicts.

\iqwrong{Trying a parity on rows and columns separately, which handles a few bits but cannot do all six at once. Each bit of the address must be controlled by a set of squares that overlaps the others in the right way, and exclusive-or gives you that.}
\iqwhy{This is a perfect code. One changed bit carries exactly $\log_2 n$ bits of information, and the same idea underlies the Hamming code used to correct errors in memory.}
\end{solution}

\begin{solution}{iq:iv1:logic-riddles:35}
Among 5 cards, two share a suit. The assistant hides one of those and puts the other first, to fix the suit and a base rank. The other three cards, in one of 6 orders, tell the distance from base to hidden. A counting argument allows up to 124 cards.

There are 13 ranks in a circle, so one of the two same-suit cards is at most 6 steps clockwise from the other\index{pigeonhole principle}. Show the card from which you count, and hide the other. Order the remaining three, using a fixed order of all cards, to encode the distance $d$ from 1 to 6, since $3! = 6$. The magician reads the first card, adds $d$ on the circle, and keeps the suit. A test with 40\,000 random hands recovers the hidden card every time.

For the counting bound with a deck of $n$ cards, each hand must be sent to a different row of 4 cards. There are $\binom{n}{5}$ hands, and $n(n - 1)(n - 2)(n - 3)$ rows. That needs $\frac{n - 4}{120} \le 1$, and so $n \le 124$. It can be reached: each hand can use 120 rows (5 choices of a hidden card times 4! orders), each row suits $n - 4$ hands, and with $n - 4 \le 120$, a matching exists by Hall's theorem\index{Hall's marriage theorem}. That proof gives no easy formula.

\iqwrong{Saying that 4! = 24 orders cannot code a card out of 48. The assistant also chooses which card to hide, and that gives 5 times the freedom. With the suit given by the first card, you only need a small part of that.}
\end{solution}
